\section{lanqiao刷题总结}
%\textcolor{blue}{\textbf{271A}}
%\begin{lstlisting}[language=C++]\end{lstlisting}
%\textcolor{red}{\textbf{注意：}}
在本节中，我将总结我在 lanqiao 上做题的经验。

\subsection{遇到的问题以及学到的方法}
\begin{itemize}
\subsubsection{最少回文子串}
\item 题目链接:\href{https://www.lanqiao.cn/problems/21192/learning/?contest_id=274}{最少子串回文数(点击此处)};方法：1、\textcolor{red}{\textbf{回文数特点(非常重要！！！):}}最多允许一个字符出现奇数次(要么所有字符的出现次数都为偶数，要么只能有一个字符出现次数为奇数(\textcolor{blue}{\textbf{出现在中心位置}}));2、总结:若所有字符出现次数均为偶数则答案为1；否则答案是出现次数为奇数的字符的个数；
	
代码示例:\\
\begin{lstlisting}[language=C++]
#include <bits/stdc++.h>
using namespace std;
int main(){
	int n;
	string s;
	cin>>n>>s;
	vector<int> v(26,0);
	for(auto x:s){
		v[x-'a']++;
	}
	int odd_cnt=0;
	for(auto ch:v){
		if(ch&1){
			odd_cnt++;
		}
	}
	if(odd_cnt==0) cout<<1<<endl;
	else cout<<odd_cnt<<endl;
	return 0;
}
\end{lstlisting}
\subsubsection{合理配置}
\item 题目链接 \href{https://www.lanqiao.cn/problems/21193/learning/?contest_id=274}{合理配置(点击此处)}；思路用总的袋子数-合理配置的数量=幽默配置的数量；
	
代码示例：\\
\newpage
\begin{lstlisting}[language=C++]
#include <bits/stdc++.h>
#define endl '\n'
using namespace std;
using ll=long long;
int main(){
	ll a,b;
	cin>>a>>b;
	ll ans=INT_MAX;
	for(ll i=0;i<=min(a,b);i++){
		ans=min(ans,abs(a+b-4*i));
	}
	cout<<ans<<endl;
	return 0;
}
\end{lstlisting}
\end{itemize}





















